\documentclass[tikz,border=4pt]{standalone}
\usepackage{ctex}
\usepackage{amsmath,amssymb}
\usepackage{esint}
\usetikzlibrary{calc,arrows.meta}
\begin{document}
\begin{tikzpicture}[
  x={(1cm,0cm)}, y={(0cm,1cm)}, z={(-0.45cm,-0.38cm)},
  line cap=round, line join=round,
  >=Stealth
]
% 坐标轴
\draw[->] (-0.5,0,0) -- (4.5,0,0) node[right] {$x$};
\draw[->] (0,-0.5,0) -- (0,4,0) node[above] {$y$};
\draw[->] (0,0,-0.5) -- (0,0,3.5) node[above left] {$z$};

% 球体轮廓（部分可见面）
\def\rr{2.0}
% 前半球（可见）
\draw[blue!40, thick]
  plot[domain=0:180, samples=30, variable=\ph]
    ({\rr*sin(\ph)},{0},{\rr*cos(\ph)});
\draw[blue!40, thick]
  plot[domain=0:180, samples=30, variable=\ph]
    ({0},{\rr*sin(\ph)},{\rr*cos(\ph)});
% 赤道
\draw[blue!40]
  plot[domain=0:360, samples=40, variable=\th]
    ({\rr*cos(\th)},{\rr*sin(\th)},0);
\draw[blue!30, dashed]
  plot[domain=0:360, samples=40, variable=\th]
    ({\rr*cos(\th)},0,{\rr*sin(\th)});

\fill[blue!10, opacity=0.3]
  plot[domain=0:360, samples=40, variable=\th]
    ({\rr*cos(\th)},{\rr*sin(\th)},0) -- cycle;

\node[blue!70!black, font=\small] at (1.2,1.2,1.5) {$V$};
\node[blue!60!black, font=\small] at (2.5,1.8,0.3) {$\partial V$（外法向）};

% 外法向箭头
\foreach \th in {0,90,180,270} {
  \pgfmathsetmacro{\nx}{cos(\th)}
  \pgfmathsetmacro{\ny}{sin(\th)}
  \draw[->, orange, thick]
    ({\rr*\nx},{\rr*\ny},0) -- ++({0.6*\nx},{0.6*\ny},0);
}
\draw[->, orange, thick] (0,0,\rr) -- ++(0,0,0.6);

% 散度（内部源）示意
\foreach \xf/\yf/\zf in {0.5/0.5/0.5, -0.5/0.5/0.3, 0.5/-0.5/0.4,
                           -0.5/-0.5/0.5, 0/0/1.0, 0/0/-0.5} {
  \draw[->, red!50, thin]
    (\xf,\yf,\zf) -- ++({0.25*\xf/1},{0.25*\yf/1},{0.25*\zf/1});
}
\node[red!80!black, font=\small] at (-1.5,0,1.5) {$\nabla\cdot\mathbf{F}$（源密度）};

% Gauss 公式
\node[font=\small, draw=gray!50, fill=gray!10, rounded corners=2pt, align=center]
  at (-1,3.8,3.2)
  {\textbf{Gauss 散度定理}\\
   $\oiint_{\partial V}\mathbf{F}\cdot d\mathbf{S} = \iiint_V \nabla\cdot\mathbf{F}\,dV$\\
   $\partial V$ 取外法向；$\mathbf{F}$ 须在 $V$ 上 $C^1$};
\end{tikzpicture}
\end{document}
